% Pass 3 proposal, Part A: argument units A01--A08.
% This is a proposal fragment. Canonical IDs, labels, and editorial ledgers
% remain the responsibility of the chapter editor.

% YIN-SOURCE: id=PROPOSAL-A-U001; arguments=YIN253A-C02-A01; notes=253a:10-11; pdf=20-22; video=96lN2omwit4:00:01:26.700-00:23:41.120; transcript=YIN253A-C02-T000003..YIN253A-C02-T000039; equations=YIN253A-C02-EQ001,YIN253A-C02-EQ002; class=SOURCE_COMPOSITE
\subsection{From Hamiltonian evolution to the transition kernel}

% YIN-SOURCE: id=PROPOSAL-A-U002; arguments=YIN253A-C02-A01; notes=253a:10; pdf=21; video=96lN2omwit4:00:02:35.879-00:08:14.060; transcript=YIN253A-C02-T000005..YIN253A-C02-T000013; equations=YIN253A-C02-EQ001; class=SOURCE_COMPOSITE
Classical mechanics has two equivalent descriptions.  In the Lagrangian
description the configuration and velocity, $q$ and $\dot q$, specify a state,
and the equations of motion follow by extremizing the action.  In the
Hamiltonian description the state is specified by $q$ and its conjugate
momentum $p$, while the Hamiltonian generates time evolution through the
Poisson bracket.  Hopefully this is not unfamiliar from classical mechanics.

% YIN-SOURCE: id=PROPOSAL-A-U003; arguments=YIN253A-C02-A01; notes=253a:10; pdf=21; video=96lN2omwit4:00:02:35.879-00:07:55.139; transcript=YIN253A-C02-T000005..YIN253A-C02-T000012; equations=YIN253A-C02-EQ001; class=SOURCE_COMPOSITE
\begin{align}
  S[q] &= \int dt\,L(q,\dot q),
  & \dot f &= \{H,f\}_{\mathrm{Poisson}},
  \label{eq:proposal-classical-descriptions}\\
  \dot q &= \frac{\partial H}{\partial p},
  & \dot p &= -\frac{\partial H}{\partial q}.
  \label{eq:proposal-hamilton-equations}
\end{align}

% YIN-SOURCE: id=PROPOSAL-A-U004; arguments=YIN253A-C02-A01; notes=253a:10; pdf=21; video=96lN2omwit4:00:06:55.139-00:08:14.060; transcript=YIN253A-C02-T000011..YIN253A-C02-T000013; equations=YIN253A-C02-EQ001; class=SOURCE_COMPOSITE
The bridge is the Legendre transform.  With the momentum fixed by
$p=\partial L/\partial\dot q$, the expression $p\dot q-L$ depends on $p$ and
$q$.  Differentiating it gives Hamilton's equations once the Euler--Lagrange
equation is used.

% YIN-SOURCE: id=PROPOSAL-A-U005; arguments=YIN253A-C02-A01; notes=253a:10; pdf=21; video=96lN2omwit4:00:06:55.139-00:08:14.060; transcript=YIN253A-C02-T000011..YIN253A-C02-T000013; equations=YIN253A-C02-EQ001; class=SOURCE_COMPOSITE
\begin{equation}
  p=\frac{\partial L}{\partial\dot q},
  \qquad
  H(p,q)=
  \left.p\dot q-L(q,\dot q)\right|_{\frac{\partial L}{\partial\dot q}=p}.
  \label{eq:proposal-legendre-transform}
\end{equation}

% YIN-SOURCE: id=PROPOSAL-A-U006; arguments=YIN253A-C02-A01; notes=253a:10-11; pdf=21-22; video=96lN2omwit4:00:10:19.380-00:15:00.420; transcript=YIN253A-C02-T000019..YIN253A-C02-T000028; equations=YIN253A-C02-EQ001,YIN253A-C02-EQ002; class=SOURCE_COMPOSITE
Quantization begins with a Hilbert space and an operator $\hat H$ obtained from
the classical Hamiltonian after an ordering prescription has been chosen.  Of
course, there are two different ways to describe time evolution.  States
evolve in the Schr\"odinger picture.  The Heisenberg picture keeps the states
fixed and evolves the observables by conjugation.

% YIN-SOURCE: id=PROPOSAL-A-U007; arguments=YIN253A-C02-A01; notes=253a:10-11; pdf=21-22; video=96lN2omwit4:00:10:44.579-00:15:00.420; transcript=YIN253A-C02-T000021..YIN253A-C02-T000028; equations=YIN253A-C02-EQ001,YIN253A-C02-EQ002; class=SOURCE_COMPOSITE
\begin{align}
  &\ket{\psi}\in\mathcal H,
  \qquad H(p,q)\longrightarrow\hat H,
  \label{eq:proposal-quantum-bridge}\\
  &\ket{\psi}\longmapsto U(t)\ket{\psi},
  \qquad
  U(t)=e^{-\frac{i}{\hbar}\hat Ht},
  \label{eq:proposal-schrodinger-evolution}\\
  &O(t)=U(t)^{-1}OU(t),
  \qquad
  \dot O(t)=\frac{i}{\hbar}[\hat H,O(t)].
  \label{eq:proposal-heisenberg-evolution}
\end{align}

% YIN-SOURCE: id=PROPOSAL-A-U008; arguments=YIN253A-C02-A01; notes=253a:11; pdf=22; video=96lN2omwit4:00:14:59.220-00:19:39.299; transcript=YIN253A-C02-T000029..YIN253A-C02-T000034; equations=YIN253A-C02-EQ002; class=SOURCE_COMPOSITE
The Lagrangian formulation of quantum mechanics requires another route, since
there is no operator that corresponds to the Lagrangian.  Consider, for
example, a particle with
$\hat H=\hat p^{\,2}/(2m)+V(\hat q)$.  The coordinate eigenstates form a
complete basis.  Therefore the matrix elements of $U(T)$ determine its action
on every state.

% YIN-SOURCE: id=PROPOSAL-A-U009; arguments=YIN253A-C02-A01; notes=253a:11; pdf=22; video=96lN2omwit4:00:18:53.220-00:23:41.120; transcript=YIN253A-C02-T000034..YIN253A-C02-T000039; equations=YIN253A-C02-EQ002; class=SOURCE_COMPOSITE
The formal Lagrangian expression for this kernel is a sum over paths with fixed
endpoints.  Its phase is fixed by the classical action.  At this stage the
symbol $[Dq]$ is only an instruction whose meaning remains to be constructed.

% YIN-SOURCE: id=PROPOSAL-A-U010; arguments=YIN253A-C02-A01; notes=253a:11; pdf=22; video=96lN2omwit4:00:18:53.220-00:23:41.120; transcript=YIN253A-C02-T000034..YIN253A-C02-T000039; equations=YIN253A-C02-EQ002; class=SOURCE_COMPOSITE
\begin{equation}
  \bra{q_f}U(T)\ket{q_i}
  =\int_{q(0)=q_i}^{q(T)=q_f}[Dq]\,
  e^{\frac{i}{\hbar}S[q]},
  \qquad
  S[q]=\int_0^Tdt\,L(q,\dot q).
  \label{eq:proposal-formal-kernel}
\end{equation}

% SOURCE-CONFLICT: EQ002 and note p.11 deliberately leave [Dq] undefined.

% YIN-SOURCE: id=PROPOSAL-A-U011; arguments=YIN253A-C02-A02; notes=253a:12-14; pdf=23-25; video=96lN2omwit4:00:23:53.000-01:20:19.219; transcript=YIN253A-C02-T000040..YIN253A-C02-T000139; equations=YIN253A-C02-EQ003,YIN253A-C02-EQ004; class=SOURCE_COMPOSITE
\subsection{Time slicing and the regulated measure}

% YIN-SOURCE: id=PROPOSAL-A-U012; arguments=YIN253A-C02-A02; notes=253a:12-13; pdf=23-24; video=96lN2omwit4:00:23:53.000-00:29:25.520; transcript=YIN253A-C02-T000040..YIN253A-C02-T000050; equations=YIN253A-C02-EQ003,YIN253A-C02-EQ004; class=SOURCE_COMPOSITE
There is an elephant in the room: what the heck is this thing here?  The
functional measure is defined by a regulator.  Divide the interval into $N$
steps of length $T/N$, insert a coordinate resolution of the identity at every
intermediate time, and keep the resulting finite-dimensional integral until
the calculation has been performed.

% YIN-SOURCE: id=PROPOSAL-A-U013; arguments=YIN253A-C02-A02; notes=253a:12; pdf=23; video=96lN2omwit4:00:29:45.000-00:33:13.220; transcript=YIN253A-C02-T000051..YIN253A-C02-T000056; equations=YIN253A-C02-EQ003; class=SOURCE_COMPOSITE
\begin{align}
  \bra{q_f}U(T)\ket{q_i}
  &={}
  \int\prod_{n=1}^{N-1}d^Dq_n
  \prod_{n=0}^{N-1}
  \bra{q_{n+1}}
  e^{-\frac{i}{\hbar}\hat H\frac{T}{N}}
  \ket{q_n},
  \label{eq:proposal-coordinate-slicing}\\
  q_f&\equiv q_N,
  \qquad q_i\equiv q_0.
\end{align}

% YIN-SOURCE: id=PROPOSAL-A-U014; arguments=YIN253A-C02-A02; notes=253a:12; pdf=23; video=96lN2omwit4:00:33:27.299-00:35:58.160; transcript=YIN253A-C02-T000057..YIN253A-C02-T000060; equations=YIN253A-C02-EQ003; class=SOURCE_COMPOSITE
Insert a momentum basis into every short-time factor.  With the normalization
used here, the momentum resolution and the coordinate-space wavefunction of a
momentum eigenstate are

% YIN-SOURCE: id=PROPOSAL-A-U015; arguments=YIN253A-C02-A02; notes=253a:12; pdf=23; video=96lN2omwit4:00:33:27.299-00:35:58.160; transcript=YIN253A-C02-T000057..YIN253A-C02-T000060; equations=YIN253A-C02-EQ003; class=SOURCE_COMPOSITE
\begin{equation}
  \mathbf 1
  =\int\frac{d^Dp_n}{(2\pi\hbar)^D}\ket{p_n}\bra{p_n},
  \qquad
  \braket{q_{n+1}}{p_n}
  =e^{\frac{i}{\hbar}p_n\cdot q_{n+1}}.
  \label{eq:proposal-momentum-resolution}
\end{equation}

% YIN-SOURCE: id=PROPOSAL-A-U016; arguments=YIN253A-C02-A02; notes=253a:12; pdf=23; video=96lN2omwit4:00:36:00.660-00:43:24.260; transcript=YIN253A-C02-T000061..YIN253A-C02-T000073; equations=YIN253A-C02-EQ003; class=SOURCE_COMPOSITE
For a short interval one expands the evolution operator and places the
$\hat p$'s and $\hat q$'s in the ordering used to define the ordinary function
$H(q_n,p_n)$.  Multiplying all short-time factors turns their phases into a
sum.  The regulated phase-space kernel is then

% YIN-SOURCE: id=PROPOSAL-A-U017; arguments=YIN253A-C02-A02; notes=253a:12; pdf=23; video=96lN2omwit4:00:36:00.660-00:46:53.240; transcript=YIN253A-C02-T000061..YIN253A-C02-T000079; equations=YIN253A-C02-EQ003; class=SOURCE_COMPOSITE
\begin{align}
  \bra{q_f}U(T)\ket{q_i}
  ={}&\int\prod_{n=1}^{N-1}d^Dq_n
  \prod_{n=0}^{N-1}\frac{d^Dp_n}{(2\pi\hbar)^D}
  \nonumber\\[-2mm]
  &\times
  \exp\!\left\{
  \frac{i}{\hbar}\sum_{n=0}^{N-1}
  \left[
    p_n\cdot(q_{n+1}-q_n)
    -H(q_n,p_n)\frac{T}{N}
    +O\!\left(\hbar,\left(\frac{T}{N}\right)^2\right)
  \right]\right\}.
  \label{eq:proposal-finite-phase-space-kernel}
\end{align}

% SOURCE-CONFLICT: p.23 places O(hbar,(T/N)^2) inside the finite-N sum. The
% hbar-shaped glyph remains source-local, and no different order symbol is
% inferred here.

% YIN-SOURCE: id=PROPOSAL-A-U018; arguments=YIN253A-C02-A02; notes=253a:12-13; pdf=23-24; video=96lN2omwit4:00:39:18.260-00:56:57.740; transcript=YIN253A-C02-T000067..YIN253A-C02-T000097; equations=YIN253A-C02-EQ003,YIN253A-C02-EQ004; class=SOURCE_COMPOSITE
There is no such sacred ordering.  An alternative placement of momentum and
coordinate resolutions gives another symbol, $H'$, for the same operator.
The two symbols can differ at order $\hbar$.  Terms which look small at one
lattice step must therefore remain until their effect on the full product has
been determined.  You drop this term, you're changing the answer.

% YIN-SOURCE: id=PROPOSAL-A-U019; arguments=YIN253A-C02-A02; notes=253a:13; pdf=24; video=96lN2omwit4:00:53:02.520-00:56:57.740; transcript=YIN253A-C02-T000090..YIN253A-C02-T000097; equations=YIN253A-C02-EQ004; class=SOURCE_COMPOSITE
\begin{align}
  \bra{p_n}\hat H\ket{q_n}
  &=H(p_n,q_n)\braket{p_n}{q_n},
  \label{eq:proposal-left-ordering}\\
  \bra{q_n}\hat H\ket{p_n}
  &=H'(p_n,q_n)\braket{q_n}{p_n},
  \qquad H'-H=O(\hbar).
  \label{eq:proposal-right-ordering}
\end{align}

% SOURCE-CONFLICT: p.23 uses H(q_n,p_n), while p.24 uses H(p_n,q_n) and
% H'(p_n,q_n). Each argument order is retained in its source-local role.

% YIN-SOURCE: id=PROPOSAL-A-U020; arguments=YIN253A-C02-A02; notes=253a:13; pdf=24; video=96lN2omwit4:00:46:55.680-00:58:18.680; transcript=YIN253A-C02-T000080..YIN253A-C02-T000100; equations=YIN253A-C02-EQ004; class=SOURCE_COMPOSITE
The continuum notation records this limiting prescription.  It does not
introduce a separate infinite-dimensional measure.  A change of lattice
scheme may require an order-$\hbar$ counterterm in $H$ so that the regulated
limits describe the same quantum system.

% YIN-SOURCE: id=PROPOSAL-A-U021; arguments=YIN253A-C02-A02; notes=253a:13; pdf=24; video=96lN2omwit4:00:46:55.680-00:58:18.680; transcript=YIN253A-C02-T000080..YIN253A-C02-T000100; equations=YIN253A-C02-EQ004; class=SOURCE_COMPOSITE
\begin{equation}
  \bra{q_f}U(T)\ket{q_i}
  =\int\mathcal Dq\,\mathcal Dp\,
  \exp\!\left\{
    \frac{i}{\hbar}\int_0^Tdt\,
    \bigl(p\cdot\dot q-H(q,p)\bigr)
  \right\},
  \label{eq:proposal-continuum-phase-space}
\end{equation}

% YIN-SOURCE: id=PROPOSAL-A-U022; arguments=YIN253A-C02-A02; notes=253a:14; pdf=25; video=96lN2omwit4:01:04:32.819-01:11:21.260; transcript=YIN253A-C02-T000113..YIN253A-C02-T000124; equations=YIN253A-C02-EQ004; class=SOURCE_COMPOSITE
When the Hamiltonian is quadratic in $p$, every momentum integral is Gaussian.
For $H=p^2/(2m)+V(q)$, completing the square isolates a standard oscillatory
Gaussian and leaves the Legendre transform in the exponent.

% YIN-SOURCE: id=PROPOSAL-A-U023; arguments=YIN253A-C02-A02; notes=253a:14; pdf=25; video=96lN2omwit4:01:05:09.480-01:11:21.260; transcript=YIN253A-C02-T000114..YIN253A-C02-T000124; equations=YIN253A-C02-EQ004; class=SOURCE_COMPOSITE
\begin{align}
  p\dot q-\frac{p^2}{2m}-V(q)
  &=-\frac{(p-m\dot q)^2}{2m}
    +\frac{m}{2}\dot q^{\,2}-V(q),
  \label{eq:proposal-complete-momentum-square}\\
  \int\mathcal Dq\,\mathcal Dp\,
  e^{\frac{i}{\hbar}\int dt\,(p\dot q-H)}
  &=\int[Dq]\,
  e^{\frac{i}{\hbar}\int dt\,L(q,\dot q)}.
  \label{eq:proposal-momentum-reduction}
\end{align}

% YIN-SOURCE: id=PROPOSAL-A-U024; arguments=YIN253A-C02-A02; notes=253a:14; pdf=25; video=96lN2omwit4:01:07:19.619-01:10:25.880; transcript=YIN253A-C02-T000118..YIN253A-C02-T000122; equations=YIN253A-C02-EQ004; class=SOURCE_COMPOSITE
In generalized coordinates the kinetic metric depends on position.  Momentum
integration then contributes a determinant to the configuration-space
measure.  This is the familiar volume element associated with
$G_{\alpha\beta}(\xi)$, appearing here at every regulated time slice.

% YIN-SOURCE: id=PROPOSAL-A-U025; arguments=YIN253A-C02-A02; notes=253a:14; pdf=25; video=96lN2omwit4:01:07:19.619-01:10:25.880; transcript=YIN253A-C02-T000118..YIN253A-C02-T000122; equations=YIN253A-C02-EQ004; class=SOURCE_COMPOSITE
\begin{align}
  L(\xi,\dot\xi)
  &=\frac12G_{\alpha\beta}(\xi)
    \dot\xi^\alpha\dot\xi^\beta-V(\xi),
  \label{eq:proposal-generalized-lagrangian}\\
  H(p,\xi)
  &=\frac12\bigl(G^{-1}(\xi)\bigr)^{\alpha\beta}
    p_\alpha p_\beta+V(\xi)+O(\hbar),
  \label{eq:proposal-generalized-hamiltonian}\\
  [D\xi]
  &\leadsto
  \prod_{n=1}^{\widetilde N}
  \frac{d\xi_n^\alpha}
  {\sqrt{2\pi i\hbar\,\frac{T}{N}}}
  \sqrt{\det G_{\alpha\beta}(\xi_n)}\,.
  \label{eq:proposal-generalized-measure}
\end{align}

% SOURCE-CONFLICT: p.25 writes a product through \widetilde N while retaining
% T/N in the normalization. It also alternates P(t) in the determinant note
% with p in the phase-space formulas. These page-local marks are withheld
% from the prose.

% YIN-SOURCE: id=PROPOSAL-A-U026; arguments=YIN253A-C02-A03; notes=253a:15-16; pdf=26-27; video=uzixOflp0tY:00:02:50.879-00:38:28.460; transcript=YIN253A-C02-T000158..YIN253A-C02-T000218; equations=YIN253A-C02-EQ005; class=SOURCE_COMPOSITE
\subsection{Spectral correlators and Euclidean time}

% YIN-SOURCE: id=PROPOSAL-A-U027; arguments=YIN253A-C02-A03; notes=253a:15; pdf=26; video=uzixOflp0tY:00:02:50.879-00:10:52.579; transcript=YIN253A-C02-T000158..YIN253A-C02-T000170; equations=YIN253A-C02-EQ005; class=SOURCE_COMPOSITE
We next use the formalism to study observables that capture physical quantities
of interest, such as the energy spectrum of a quantum-mechanical system.  A
particularly useful observable is the ground-state two-point function.  Insert
a complete basis of energy eigenstates between its two operators.

% YIN-SOURCE: id=PROPOSAL-A-U028; arguments=YIN253A-C02-A03; notes=253a:15; pdf=26; video=uzixOflp0tY:00:04:05.159-00:10:52.579; transcript=YIN253A-C02-T000160..YIN253A-C02-T000170; equations=YIN253A-C02-EQ005; class=SOURCE_COMPOSITE
\begin{align}
  \hat H\ket{n}&=E_n\ket{n},
  \label{eq:proposal-energy-eigenstates}\\
  G(t)
  &\equiv\bra{0}\hat q_\beta(t)\hat q_\beta(0)\ket{0}
  \nonumber\\
  &=\sum_n
  \bra{0}\hat q_\beta(0)\ket{n}
  \bra{n}\hat q_\beta(0)\ket{0}\,
  e^{-\frac{i}{\hbar}(E_n-E_0)t}.
  \label{eq:proposal-spectral-correlator}
\end{align}

% YIN-SOURCE: id=PROPOSAL-A-U029; arguments=YIN253A-C02-A03; notes=253a:15; pdf=26; video=uzixOflp0tY:00:09:55.080-00:11:08.420; transcript=YIN253A-C02-T000168..YIN253A-C02-T000171; equations=YIN253A-C02-EQ005; class=SOURCE_COMPOSITE
Every state with a nonzero matrix element contributes an oscillation whose
frequency is the gap $E_n-E_0$.  For the harmonic oscillator,
$\hat q\ket{0}$ overlaps only with the first excited state.  A generic system
can display many gaps in the same correlator.

% YIN-SOURCE: id=PROPOSAL-A-U030; arguments=YIN253A-C02-A03; notes=253a:15-16; pdf=26-27; video=uzixOflp0tY:00:11:09.000-00:18:32.720; transcript=YIN253A-C02-T000172..YIN253A-C02-T000182; equations=YIN253A-C02-EQ005; class=SOURCE_COMPOSITE
Each term can be evaluated at complex $t$.  There is nothing wrong about that.
But there is a question of when the infinite sum will converge.  Since
$E_n-E_0\geq0$, continuation through the lower half-plane,
$\operatorname{Im}t<0$, suppresses the high-energy terms.  On the negative
imaginary axis one writes $t=-i\tau$ with $\tau>0$.

% YIN-SOURCE: id=PROPOSAL-A-U031; arguments=YIN253A-C02-A03; notes=253a:15-16; pdf=26-27; video=uzixOflp0tY:00:12:43.740-00:18:32.720; transcript=YIN253A-C02-T000175..YIN253A-C02-T000182; equations=YIN253A-C02-EQ005; class=SOURCE_COMPOSITE
\begin{equation}
  G(-i\tau)
  =\sum_n
  \bra{0}\hat q(0)\ket{n}
  \bra{n}\hat q(0)\ket{0}\,
  e^{-\frac{1}{\hbar}(E_n-E_0)\tau},
  \qquad \tau>0.
  \label{eq:proposal-euclidean-spectral-correlator}
\end{equation}

% SOURCE-CONFLICT: pp.26-27 alternate unbarred and barred energy symbols.
% The proposal uses E_n consistently while preserving the local accent change
% in this comment.

% YIN-SOURCE: id=PROPOSAL-A-U032; arguments=YIN253A-C02-A03; notes=253a:16; pdf=27; video=uzixOflp0tY:00:24:42.659-00:28:27.260; transcript=YIN253A-C02-T000194..YIN253A-C02-T000200; equations=YIN253A-C02-EQ005; class=SOURCE_COMPOSITE
Euclidean evolution also provides a direct projection onto the ground state.
Expand a generic state as $\ket{\psi}=\sum_n c_n\ket{n}$, with $c_0\neq0$.
Every excited component receives an exponentially small factor as
$T\to\infty$.

% YIN-SOURCE: id=PROPOSAL-A-U033; arguments=YIN253A-C02-A03; notes=253a:16; pdf=27; video=uzixOflp0tY:00:24:42.659-00:28:27.260; transcript=YIN253A-C02-T000194..YIN253A-C02-T000200; equations=YIN253A-C02-EQ005; class=SOURCE_COMPOSITE
\begin{equation}
  e^{-\frac{1}{\hbar}(\hat H-E_0)T}\ket{\psi}
  =c_0\ket{0}+
  \sum_{n>0}c_n
  e^{-\frac{1}{\hbar}(E_n-E_0)T}\ket{n}
  \xrightarrow[T\to\infty]{}c_0\ket{0}.
  \label{eq:proposal-ground-state-projection}
\end{equation}

% YIN-SOURCE: id=PROPOSAL-A-U034; arguments=YIN253A-C02-A03; notes=253a:16; pdf=27; video=uzixOflp0tY:00:25:26.960-00:26:29.539; transcript=YIN253A-C02-T000195..YIN253A-C02-T000196; equations=YIN253A-C02-EQ005; class=SOURCE_COMPOSITE
This operator is going to kill everything but the ground state.  Everything
that is not the ground state will be exponentially suppressed.  The overlap
condition $c_0\neq0$ is the only input needed for this projection.

% YIN-SOURCE: id=PROPOSAL-A-U035; arguments=YIN253A-C02-A04; notes=253a:17-18; pdf=28-29; video=uzixOflp0tY:00:38:28.680-01:07:17.000; transcript=YIN253A-C02-T000219..YIN253A-C02-T000282; equations=YIN253A-C02-EQ006; class=SOURCE_COMPOSITE
\subsection{The normalized Euclidean path integral}

% YIN-SOURCE: id=PROPOSAL-A-U036; arguments=YIN253A-C02-A04; notes=253a:17; pdf=28; video=uzixOflp0tY:00:38:28.680-00:47:03.480; transcript=YIN253A-C02-T000219..YIN253A-C02-T000238; equations=YIN253A-C02-EQ006; class=SOURCE_COMPOSITE
Place the two operator insertions between long Euclidean evolution factors and
divide by the same expression without insertions.  The common ground-state
overlap and its normalization cancel.  In the operator expression the
insertions carry hats.

% YIN-SOURCE: id=PROPOSAL-A-U037; arguments=YIN253A-C02-A04; notes=253a:16-17; pdf=27-28; video=uzixOflp0tY:00:26:46.320-00:46:50.300; transcript=YIN253A-C02-T000197..YIN253A-C02-T000237; equations=YIN253A-C02-EQ006; class=SOURCE_COMPOSITE
\begin{equation}
  G_E(\tau)
  =\lim_{T\to\infty}
  \frac{
  \bra{\psi}
  e^{-\frac{1}{\hbar}\hat H(T-\tau)}
  \hat q(0)e^{-\frac{1}{\hbar}\hat H\tau}
  \hat q(0)e^{-\frac{1}{\hbar}\hat HT}
  \ket{\psi}}
  {\bra{\psi}e^{-\frac{2}{\hbar}\hat HT}\ket{\psi}}.
  \label{eq:proposal-normalized-operator-correlator}
\end{equation}

% SOURCE-CONFLICT: note p.17 has a positive first exponential in its top
% operator line and writes the second insertion without a hat in the long-T
% ratio. The normalized equation index and the operator/path-variable
% distinction govern the proposal; the page-local marks remain here.

% YIN-SOURCE: id=PROPOSAL-A-U038; arguments=YIN253A-C02-A04; notes=253a:17; pdf=28; video=uzixOflp0tY:00:38:28.680-00:47:03.480; transcript=YIN253A-C02-T000219..YIN253A-C02-T000238; equations=YIN253A-C02-EQ006; class=SOURCE_COMPOSITE
Time slicing converts an insertion of $\hat q$ into multiplication by the
coordinate variable at that slice.  Thus the hats disappear inside the
functional integral.  The whole point is that analytic continuation is
unambiguous if you have the function for imaginary time.

% YIN-SOURCE: id=PROPOSAL-A-U039; arguments=YIN253A-C02-A04; notes=253a:17-18; pdf=28-29; video=uzixOflp0tY:00:38:28.680-00:57:30.800; transcript=YIN253A-C02-T000219..YIN253A-C02-T000261; equations=YIN253A-C02-EQ006; class=SOURCE_COMPOSITE
\begin{equation}
  G_E(\tau)
  =\lim_{T\to\infty}
  \frac{
    \displaystyle\int[Dq]\,
    e^{-\frac{1}{\hbar}S^E[q]}q(\tau)q(0)}
  {
    \displaystyle\int[Dq]\,
    e^{-\frac{1}{\hbar}S^E[q]}}.
  \label{eq:proposal-normalized-euclidean-path-integral}
\end{equation}

% SOURCE-CONFLICT: the Lorentzian denominator on note p.17 omits dt. The
% omission is not promoted into the normalized Euclidean expression.

% YIN-SOURCE: id=PROPOSAL-A-U040; arguments=YIN253A-C02-A04; notes=253a:17-18; pdf=28-29; video=uzixOflp0tY:00:47:04.200-00:54:42.119; transcript=YIN253A-C02-T000239..YIN253A-C02-T000255; equations=YIN253A-C02-EQ006; class=SOURCE_COMPOSITE
Under the Wick map $t=-i\tau$, the time derivative becomes
$\dot q=i\partial_\tau q$ and $dt=-i\,d\tau$.  Define the Euclidean
Lagrangian so that the oscillatory Lorentzian phase becomes a decaying weight.

% YIN-SOURCE: id=PROPOSAL-A-U041; arguments=YIN253A-C02-A04; notes=253a:17-18; pdf=28-29; video=uzixOflp0tY:00:47:04.200-00:54:42.119; transcript=YIN253A-C02-T000239..YIN253A-C02-T000255; equations=YIN253A-C02-EQ006; class=SOURCE_COMPOSITE
\begin{align}
  q(t)&\mathrel{\overset{t=-i\tau}{\rightsquigarrow}}q(\tau),
  & \dot q&=+i\partial_\tau q,
  \label{eq:proposal-wick-map}\\
  L[q,\dot q]&\equiv-L^E[q,\partial_\tau q],
  & S^E[q]&=\int d\tau\,L^E[q,\partial_\tau q],
  \label{eq:proposal-euclidean-action-definition}\\
  e^{\frac{i}{\hbar}\int L\,dt}
  &\rightsquigarrow e^{-\frac{1}{\hbar}S^E[q]}.
  \label{eq:proposal-wick-weight}
\end{align}

% SOURCE-CONFLICT: note p.18 literally writes \partial_t\dot q=
% +i\partial_\tau\dot q. The proposal uses the derivative map fixed by
% t=-i tau and preserves the page-local doubled-dot reading in this comment.

% YIN-SOURCE: id=PROPOSAL-A-U042; arguments=YIN253A-C02-A04; notes=253a:18; pdf=29; video=uzixOflp0tY:00:51:26.760-00:55:31.460; transcript=YIN253A-C02-T000249..YIN253A-C02-T000257; equations=YIN253A-C02-EQ006; class=SOURCE_COMPOSITE
For an ordinary kinetic term the Euclidean action is positive when the
potential is bounded below.  Wildly fluctuating paths carry a large action;
the integral is dominated by paths with small values of the Euclidean action.

% YIN-SOURCE: id=PROPOSAL-A-U043; arguments=YIN253A-C02-A04; notes=253a:18; pdf=29; video=uzixOflp0tY:00:51:26.760-00:55:31.460; transcript=YIN253A-C02-T000249..YIN253A-C02-T000257; equations=YIN253A-C02-EQ006; class=SOURCE_COMPOSITE
\begin{equation}
  L(q,\dot q)=\frac12\dot q^{\,2}-V(q)
  \qquad\longrightarrow\qquad
  L^E(q,\partial_\tau q)
  =\frac12(\partial_\tau q)^2+V(q).
  \label{eq:proposal-euclidean-lagrangian}
\end{equation}

% YIN-SOURCE: id=PROPOSAL-A-U044; arguments=YIN253A-C02-A04; notes=253a:18; pdf=29; video=uzixOflp0tY:01:02:11.760-01:07:17.000; transcript=YIN253A-C02-T000273..YIN253A-C02-T000282; equations=YIN253A-C02-EQ006; class=SOURCE_COMPOSITE
The remaining question concerns the measure.  A basis expansion turns a path
into a sequence of ordinary coefficients, which can be regulated by retaining
finitely many modes.  It is a nice version of quantum mechanics where some
ambiguities get resolved by the requirement of locality.  The oscillator
provides the first explicit calculation.

% YIN-SOURCE: id=PROPOSAL-A-U045; arguments=YIN253A-C02-A05; notes=253a:19-20; pdf=30-31; video=uzixOflp0tY:01:07:17.280-01:19:04.800; transcript=YIN253A-C02-T000283..YIN253A-C02-T000304; equations=YIN253A-C02-EQ007; class=SOURCE_COMPOSITE
\subsection{The oscillator as a regulated functional integral}

% YIN-SOURCE: id=PROPOSAL-A-U046; arguments=YIN253A-C02-A05; notes=253a:19; pdf=30; video=uzixOflp0tY:01:07:17.280-01:10:40.320; transcript=YIN253A-C02-T000283..YIN253A-C02-T000287; equations=YIN253A-C02-EQ007; class=SOURCE_COMPOSITE
Take the harmonic oscillator with unit mass.  Its answer is already known from
elementary quantum mechanics, but we are going to go through this painful
process using the path integral to calculate it.  Work first on a finite
Euclidean interval and impose Dirichlet boundary conditions.

% YIN-SOURCE: id=PROPOSAL-A-U047; arguments=YIN253A-C02-A05; notes=253a:19; pdf=30; video=uzixOflp0tY:01:08:15.480-01:11:59.960; transcript=YIN253A-C02-T000284..YIN253A-C02-T000290; equations=YIN253A-C02-EQ007; class=SOURCE_COMPOSITE
\begin{align}
  H&=\frac{p^2+\omega^2q^2}{2},
  &L&=\frac12\dot q^{\,2}-\frac12\omega^2q^2,
  \label{eq:proposal-oscillator-lorentzian}\\
  S^E[q]&=\int_{-T}^{T}d\tau\,
  \left[\frac12(\partial_\tau q)^2+\frac12\omega^2q^2\right],
  &q(-T)&=q(T)=0.
  \label{eq:proposal-oscillator-euclidean}
\end{align}

% YIN-SOURCE: id=PROPOSAL-A-U048; arguments=YIN253A-C02-A05; notes=253a:19-20; pdf=30-31; video=uzixOflp0tY:01:12:00.360-01:15:07.460; transcript=YIN253A-C02-T000291..YIN253A-C02-T000296; equations=YIN253A-C02-EQ007; class=SOURCE_COMPOSITE
The boundary-adapted sine functions form an orthonormal basis on $[-T,T]$.
Their normalization makes the mode coefficients ordinary Cartesian
coordinates on the regulated space of paths.

% YIN-SOURCE: id=PROPOSAL-A-U049; arguments=YIN253A-C02-A05; notes=253a:19-20; pdf=30-31; video=uzixOflp0tY:01:12:00.360-01:15:07.460; transcript=YIN253A-C02-T000291..YIN253A-C02-T000296; equations=YIN253A-C02-EQ007; class=SOURCE_COMPOSITE
\begin{align}
  q(\tau)&=\sum_{n=1}^{\infty}q_nf_n(\tau),
  \qquad
  f_n(\tau)=\frac1{\sqrt T}
  \sin\frac{n\pi(\tau+T)}{2T},
  \label{eq:proposal-dirichlet-modes}\\
  (f_n,f_m)&:=\int_{-T}^{T}d\tau\,
  f_n^*(\tau)f_m(\tau)=\delta_{nm}.
  \label{eq:proposal-dirichlet-orthonormality}
\end{align}

% YIN-SOURCE: id=PROPOSAL-A-U050; arguments=YIN253A-C02-A05; notes=253a:20; pdf=31; video=uzixOflp0tY:01:15:09.000-01:18:37.400; transcript=YIN253A-C02-T000297..YIN253A-C02-T000303; equations=YIN253A-C02-EQ007; class=SOURCE_COMPOSITE
Substitution into the action gives no mixing between different modes.  The
derivative term contributes $(n\pi/2T)^2$, while the potential contributes
$\omega^2$.  Each coefficient is therefore governed by an independent
positive Gaussian.

% YIN-SOURCE: id=PROPOSAL-A-U051; arguments=YIN253A-C02-A05; notes=253a:20; pdf=31; video=uzixOflp0tY:01:15:09.000-01:18:37.400; transcript=YIN253A-C02-T000297..YIN253A-C02-T000303; equations=YIN253A-C02-EQ007; class=SOURCE_COMPOSITE
\begin{equation}
  S^E[q]
  =\sum_{n=1}^{\infty}
  \left[
    \frac12\left(\frac{n\pi}{2T}\right)^2
    +\frac12\omega^2
  \right]q_n^2.
  \label{eq:proposal-diagonal-mode-action}
\end{equation}

% YIN-SOURCE: id=PROPOSAL-A-U052; arguments=YIN253A-C02-A05; notes=253a:18-20; pdf=29-31; video=uzixOflp0tY:01:03:36.780-01:18:37.400; transcript=YIN253A-C02-T000276..YIN253A-C02-T000303; equations=YIN253A-C02-EQ007; class=SOURCE_COMPOSITE
The product measure is the regulator.  At an intermediate stage retain the
first $M$ coefficients, perform the ordinary integral, and then take
$M\to\infty$.  Overall constants cancel in normalized expectation values.

% YIN-SOURCE: id=PROPOSAL-A-U053; arguments=YIN253A-C02-A05; notes=253a:18-20; pdf=29-31; video=uzixOflp0tY:01:03:36.780-01:18:37.400; transcript=YIN253A-C02-T000276..YIN253A-C02-T000303; equations=YIN253A-C02-EQ007; class=SOURCE_COMPOSITE
\begin{equation}
  [Dq]_M=\mathcal N_M\prod_{n=1}^{M}dq_n,
  \qquad
  \int[Dq]e^{-S^E/\hbar}(\cdots)
  :=\lim_{M\to\infty}
  \mathcal N_M\int\prod_{n=1}^{M}dq_n\,
  e^{-S^E_M/\hbar}(\cdots).
  \label{eq:proposal-mode-regulator}
\end{equation}

% SOURCE-CONFLICT: note p.18 leaves [Dq] -> product dq_n with a question
% mark. Pages 19-20 supply the explicit basis and Gaussian product. The finite
% partial product above makes that chosen regulator explicit.

% YIN-SOURCE: id=PROPOSAL-A-U054; arguments=YIN253A-C02-A06; notes=253a:21-23; pdf=32-34; video=M0py5a4RWhE:00:00:02.340-00:41:30.140; transcript=YIN253A-C02-T000321..YIN253A-C02-T000396; equations=YIN253A-C02-EQ008,YIN253A-C02-EQ009; class=SOURCE_COMPOSITE
\subsection{Mode covariance and the oscillator Green function}

% YIN-SOURCE: id=PROPOSAL-A-U055; arguments=YIN253A-C02-A06; notes=253a:21; pdf=32; video=M0py5a4RWhE:00:03:31.040-00:14:51.600; transcript=YIN253A-C02-T000325..YIN253A-C02-T000342; equations=YIN253A-C02-EQ008; class=SOURCE_COMPOSITE
Let $a_n=(n\pi/2T)^2+\omega^2$.  Odd moments vanish, and the normalized
one-mode Gaussian gives $\hbar/a_n$.  Since different $q_n$ occur in separate
factors, their covariance is diagonal.  The paths being integrated are
arbitrary functions obeying the endpoint condition.  Nowhere did I say to use
the equation of motion.

% YIN-SOURCE: id=PROPOSAL-A-U056; arguments=YIN253A-C02-A06; notes=253a:21; pdf=32; video=M0py5a4RWhE:00:05:51.000-00:14:51.600; transcript=YIN253A-C02-T000328..YIN253A-C02-T000342; equations=YIN253A-C02-EQ008; class=SOURCE_COMPOSITE
\begin{align}
  \langle q_nq_m\rangle
  &=\delta_{nm}
  \frac{
    \displaystyle\int_{-\infty}^{\infty}dq_n\,q_n^2
    e^{-\frac{a_n}{2\hbar}q_n^2}}
  {
    \displaystyle\int_{-\infty}^{\infty}dq_n\,
    e^{-\frac{a_n}{2\hbar}q_n^2}}
  \nonumber\\
  &=\delta_{nm}\,
  \frac{\hbar}{\left(\frac{n\pi}{2T}\right)^2+\omega^2}.
  \label{eq:proposal-mode-covariance}
\end{align}

% YIN-SOURCE: id=PROPOSAL-A-U057; arguments=YIN253A-C02-A06; notes=253a:21; pdf=32; video=M0py5a4RWhE:00:14:48.720-00:16:46.160; transcript=YIN253A-C02-T000343..YIN253A-C02-T000346; equations=YIN253A-C02-EQ008; class=SOURCE_COMPOSITE
Expanding both insertions in modes reduces the functional two-point function
to this covariance.  Before the large-$T$ limit it retains the dependence on
the chosen endpoints.

% YIN-SOURCE: id=PROPOSAL-A-U058; arguments=YIN253A-C02-A06; notes=253a:21; pdf=32; video=M0py5a4RWhE:00:14:48.720-00:16:46.160; transcript=YIN253A-C02-T000343..YIN253A-C02-T000346; equations=YIN253A-C02-EQ008; class=SOURCE_COMPOSITE
\begin{align}
  \langle q(\tau_1)q(\tau_2)\rangle_T
  ={}&\sum_{n=1}^{\infty}
  \frac{\hbar}{\left(\frac{n\pi}{2T}\right)^2+\omega^2}
  \frac1T
  \sin\frac{n\pi(\tau_1+T)}{2T}
  \sin\frac{n\pi(\tau_2+T)}{2T}.
  \label{eq:proposal-finite-T-covariance}
\end{align}

% YIN-SOURCE: id=PROPOSAL-A-U059; arguments=YIN253A-C02-A06; notes=253a:21-22; pdf=32-33; video=M0py5a4RWhE:00:20:27.480-00:26:25.460; transcript=YIN253A-C02-T000353..YIN253A-C02-T000364; equations=YIN253A-C02-EQ008; class=SOURCE_COMPOSITE
As $T$ grows, the allowed values $k=n\pi/(2T)$ become dense on the positive
real axis.  The sum with spacing $\pi/(2T)$ becomes an integral, while the sine
functions retain the finite-interval image terms.

% YIN-SOURCE: id=PROPOSAL-A-U060; arguments=YIN253A-C02-A06; notes=253a:21-22; pdf=32-33; video=M0py5a4RWhE:00:22:16.260-00:26:25.460; transcript=YIN253A-C02-T000356..YIN253A-C02-T000364; equations=YIN253A-C02-EQ008; class=SOURCE_COMPOSITE
\begin{align}
  k&\equiv\frac{n\pi}{2T},
  &\sum_n\frac{\pi}{2T}(\cdots)&\rightsquigarrow
  \int_0^\infty dk\,(\cdots),
  \label{eq:proposal-sum-to-integral}\\
  \langle q(\tau_1)q(\tau_2)\rangle_T
  ={}&\frac{2\hbar}{\pi}\int_0^\infty dk\,
  \frac{\sin k(\tau_1+T)\sin k(\tau_2+T)}{k^2+\omega^2}.
  \label{eq:proposal-half-line-covariance}
\end{align}

% YIN-SOURCE: id=PROPOSAL-A-U061; arguments=YIN253A-C02-A06; notes=253a:22; pdf=33; video=M0py5a4RWhE:00:24:02.220-00:28:16.100; transcript=YIN253A-C02-T000360..YIN253A-C02-T000368; equations=YIN253A-C02-EQ008,YIN253A-C02-EQ009; class=SOURCE_COMPOSITE
Writing the sines as exponentials makes the direct term and the endpoint image
term explicit.  The evenness of the denominator permits an integral over the
whole real axis.

% YIN-SOURCE: id=PROPOSAL-A-U062; arguments=YIN253A-C02-A06; notes=253a:22-23; pdf=33-34; video=M0py5a4RWhE:00:24:02.220-00:38:52.820; transcript=YIN253A-C02-T000360..YIN253A-C02-T000391; equations=YIN253A-C02-EQ008,YIN253A-C02-EQ009; class=SOURCE_COMPOSITE
\begin{align}
  \langle q(\tau_1)q(\tau_2)\rangle_T
  ={}&\frac{\hbar}{\pi}\int_{-\infty}^{\infty}
  \frac{dk}{k^2+\omega^2}
  \biggl[
  \frac14e^{ik(\tau_1-\tau_2)}
  +\frac14e^{-ik(\tau_1-\tau_2)}
  \nonumber\\[-1mm]
  &\hspace{37mm}
  -\frac14e^{ik(\tau_1+\tau_2+2T)}
  -\frac14e^{-ik(\tau_1+\tau_2+2T)}
  \biggr].
  \label{eq:proposal-full-line-covariance}
\end{align}

% YIN-SOURCE: id=PROPOSAL-A-U063; arguments=YIN253A-C02-A06; notes=253a:22-23; pdf=33-34; video=M0py5a4RWhE:00:27:59.640-00:37:13.020; transcript=YIN253A-C02-T000368..YIN253A-C02-T000389; equations=YIN253A-C02-EQ009; class=SOURCE_COMPOSITE
The elementary integral is evaluated by contour deformation.  The integral is
exactly the same provided that the function is complex analytic in between the
two contours.  For positive $\tau$, exponential decay closes the contour in
the upper half-plane and encloses $k=i\omega$.  Negative $\tau$ requires the
lower half-plane, whose clockwise orientation supplies the corresponding
sign.

% YIN-SOURCE: id=PROPOSAL-A-U064; arguments=YIN253A-C02-A06; notes=253a:22-23; pdf=33-34; video=M0py5a4RWhE:00:27:59.640-00:37:13.020; transcript=YIN253A-C02-T000368..YIN253A-C02-T000389; equations=YIN253A-C02-EQ009; class=SOURCE_COMPOSITE
\begin{align}
  \int_{-\infty}^{\infty}dk\,
  \frac{e^{ik\tau}}{k^2+\omega^2}
  &=2\pi i\,
  \underset{k\to i\omega}{\operatorname{Res}}
  \frac{e^{ik\tau}}{k^2+\omega^2}
  =\frac{\pi}{\omega}e^{-\omega\tau},
  &&\tau>0,
  \label{eq:proposal-positive-time-residue}\\
  \int_{-\infty}^{\infty}dk\,
  \frac{e^{ik\tau}}{k^2+\omega^2}
  &=\frac{\pi}{\omega}e^{\omega\tau},
  &&\tau<0.
  \label{eq:proposal-negative-time-residue}
\end{align}

% YIN-SOURCE: id=PROPOSAL-A-U065; arguments=YIN253A-C02-A06; notes=253a:23; pdf=34; video=M0py5a4RWhE:00:37:48.560-00:40:23.359; transcript=YIN253A-C02-T000390..YIN253A-C02-T000394; equations=YIN253A-C02-EQ009; class=SOURCE_COMPOSITE
Define $\tau_{12}=\tau_1-\tau_2$.  The two direct terms combine into a
function of $|\tau_{12}|$.  The image term vanishes as the endpoints recede,
and the normalized large-$T$ covariance becomes the standard Euclidean
oscillator propagator.

% YIN-SOURCE: id=PROPOSAL-A-U066; arguments=YIN253A-C02-A06; notes=253a:23; pdf=34; video=M0py5a4RWhE:00:37:48.560-00:40:23.359; transcript=YIN253A-C02-T000390..YIN253A-C02-T000394; equations=YIN253A-C02-EQ009; class=SOURCE_COMPOSITE
\begin{align}
  \langle q(\tau_1)q(\tau_2)\rangle_T
  &=\frac{\hbar}{\omega}
  \left[
    \frac12e^{-\omega|\tau_{12}|}
    -\frac12e^{-\omega(\tau_1+\tau_2+2T)}
  \right],
  \label{eq:proposal-finite-T-result}\\
  \lim_{T\to\infty}
  \langle q(\tau_1)q(\tau_2)\rangle_T
  &=\frac{\hbar}{2\omega}e^{-\omega|\tau_{12}|}.
  \label{eq:proposal-oscillator-propagator}
\end{align}

% YIN-SOURCE: id=PROPOSAL-A-U067; arguments=YIN253A-C02-A06; notes=253a:23; pdf=34; video=M0py5a4RWhE:00:38:55.020-00:40:23.359; transcript=YIN253A-C02-T000392..YIN253A-C02-T000394; equations=YIN253A-C02-EQ009; class=SOURCE_COMPOSITE
If you get confused, you can always go back to the representation in terms of
energy eigenstates and understand the meaning of the quantity being computed.
The single exponential carries the harmonic-oscillator gap $E_1-E_0=\hbar
\omega$.  A more general operator or potential produces a sum of exponentials.

% YIN-SOURCE: id=PROPOSAL-A-U068; arguments=YIN253A-C02-A07; notes=253a:24-27; pdf=35-38; video=M0py5a4RWhE:00:41:31.940-01:23:04.739; transcript=YIN253A-C02-T000397..YIN253A-C02-T000455; equations=YIN253A-C02-EQ010,YIN253A-C02-EQ011; class=SOURCE_COMPOSITE
\subsection{Gaussian identities and Wick's rule}

% YIN-SOURCE: id=PROPOSAL-A-U069; arguments=YIN253A-C02-A07; notes=253a:24; pdf=35; video=M0py5a4RWhE:00:41:31.940-00:45:54.500; transcript=YIN253A-C02-T000397..YIN253A-C02-T000405; equations=YIN253A-C02-EQ010; class=SOURCE_COMPOSITE
The mode calculation worked, but it was a pain.  We can do this once and for
all.  Let $A$ be a symmetric matrix for which the Gaussian converges, and add a
source $\mathbf J$.  Completing the square reduces the integral to the
source-free normalization.

% YIN-SOURCE: id=PROPOSAL-A-U070; arguments=YIN253A-C02-A07; notes=253a:24; pdf=35; video=M0py5a4RWhE:00:42:44.400-00:45:54.500; transcript=YIN253A-C02-T000399..YIN253A-C02-T000405; equations=YIN253A-C02-EQ010; class=SOURCE_COMPOSITE
\begin{align}
  -\frac12\mathbf q^TA\mathbf q+\mathbf J\cdot\mathbf q
  &=-\frac12(\mathbf q-A^{-1}\mathbf J)^T
  A(\mathbf q-A^{-1}\mathbf J)
  +\frac12\mathbf J^TA^{-1}\mathbf J,
  \label{eq:proposal-finite-gaussian-square}\\
  Z[\mathbf J]
  &:=\int d^N\mathbf q\,
  e^{-\frac12\mathbf q^TA\mathbf q+\mathbf J\cdot\mathbf q}
  \nonumber\\
  &=(2\pi)^{N/2}(\det A)^{-1/2}
  e^{\frac12\mathbf J^TA^{-1}\mathbf J}.
  \label{eq:proposal-finite-gaussian-generating-function}
\end{align}

% YIN-SOURCE: id=PROPOSAL-A-U071; arguments=YIN253A-C02-A07; notes=253a:24; pdf=35; video=M0py5a4RWhE:00:46:24.720-00:50:51.900; transcript=YIN253A-C02-T000406..YIN253A-C02-T000411; equations=YIN253A-C02-EQ010; class=SOURCE_COMPOSITE
Every derivative with respect to $J_n$ brings down one factor of $q_n$.
Consequently a polynomial insertion can be replaced by the corresponding
differential operator acting on $Z[\mathbf J]$.  Odd moments vanish because
the source-free measure is even.

% YIN-SOURCE: id=PROPOSAL-A-U072; arguments=YIN253A-C02-A07; notes=253a:24; pdf=35; video=M0py5a4RWhE:00:46:24.720-00:52:35.900; transcript=YIN253A-C02-T000406..YIN253A-C02-T000413; equations=YIN253A-C02-EQ010; class=SOURCE_COMPOSITE
\begin{align}
  \langle F(\mathbf q)\rangle
  &:=\frac1{Z[0]}\int d^N\mathbf q\,
  e^{-\frac12\mathbf q^TA\mathbf q}F(\mathbf q)
  \nonumber\\
  &=\left.
  \frac1{Z[0]}
  F\!\left(\frac{\partial}{\partial\mathbf J}\right)
  Z[\mathbf J]\right|_{\mathbf J=0},
  \label{eq:proposal-gaussian-source-derivatives}\\
  \langle q_n\rangle&=0,
  \qquad
  \langle q_nq_m\rangle=(A^{-1})_{nm}.
  \label{eq:proposal-finite-gaussian-two-point}
\end{align}

% YIN-SOURCE: id=PROPOSAL-A-U073; arguments=YIN253A-C02-A07; notes=253a:24-25; pdf=35-36; video=M0py5a4RWhE:00:52:45.540-00:59:44.540; transcript=YIN253A-C02-T000414..YIN253A-C02-T000420; equations=YIN253A-C02-EQ010; class=SOURCE_COMPOSITE
For four insertions there are $4!=24$ assignments of derivatives to source
factors.  Grouping the eight identical assignments for each pairing leaves
three inequivalent contractions.

% YIN-SOURCE: id=PROPOSAL-A-U074; arguments=YIN253A-C02-A07; notes=253a:24-25; pdf=35-36; video=M0py5a4RWhE:00:52:45.540-00:59:44.540; transcript=YIN253A-C02-T000414..YIN253A-C02-T000420; equations=YIN253A-C02-EQ010; class=SOURCE_COMPOSITE
\begin{align}
  \langle q_nq_mq_rq_s\rangle
  ={}&(A^{-1})_{nm}(A^{-1})_{rs}
  +(A^{-1})_{nr}(A^{-1})_{ms}
  \nonumber\\
  &+(A^{-1})_{ns}(A^{-1})_{mr}.
  \label{eq:proposal-four-point-gaussian}
\end{align}

% YIN-SOURCE: id=PROPOSAL-A-U075; arguments=YIN253A-C02-A07; notes=253a:25; pdf=36; video=M0py5a4RWhE:00:58:51.839-01:06:19.880; transcript=YIN253A-C02-T000420..YIN253A-C02-T000428; equations=YIN253A-C02-EQ010; class=SOURCE_COMPOSITE
The general rule is Wick's rule.  Pair every insertion once, replace each pair
by its two-point covariance, and sum over inequivalent pairings.  An odd number
of insertions has no complete pairing and hence gives zero.

% YIN-SOURCE: id=PROPOSAL-A-U076; arguments=YIN253A-C02-A07; notes=253a:25; pdf=36; video=M0py5a4RWhE:00:58:51.839-01:06:19.880; transcript=YIN253A-C02-T000420..YIN253A-C02-T000428; equations=YIN253A-C02-EQ010; class=SOURCE_COMPOSITE
\begin{equation}
  \langle q_{i_1}\cdots q_{i_{2r}}\rangle
  =\sum_{\text{inequivalent pairings }P}
  \prod_{(a,b)\in P}(A^{-1})_{i_ai_b},
  \qquad
  \overset{\frown}{q_nq_m}:=(A^{-1})_{nm}.
  \label{eq:proposal-wick-rule}
\end{equation}

% YIN-SOURCE: id=PROPOSAL-A-U077; arguments=YIN253A-C02-A07; notes=253a:26; pdf=37; video=M0py5a4RWhE:01:07:04.140-01:12:44.780; transcript=YIN253A-C02-T000430..YIN253A-C02-T000437; equations=YIN253A-C02-EQ011; class=SOURCE_COMPOSITE
The same algebra applies when the vector space is replaced by a linear space of
real functions.  At the level of mathematical rigor used here, the precise
function space and its suitable measure will remain implicit.  For the
oscillator, the matrix $A$ becomes a differential operator.

% YIN-SOURCE: id=PROPOSAL-A-U078; arguments=YIN253A-C02-A07; notes=253a:26; pdf=37; video=M0py5a4RWhE:01:07:04.140-01:14:38.880; transcript=YIN253A-C02-T000430..YIN253A-C02-T000440; equations=YIN253A-C02-EQ011; class=SOURCE_COMPOSITE
\begin{align}
  A\cdot g(\tau)
  &:=(-\partial_\tau^2+\omega^2)g(\tau),
  \label{eq:proposal-functional-A}\\
  A^{-1}\cdot g(\tau)
  &:=\int d\tau'\,G(\tau,\tau')g(\tau'),
  \label{eq:proposal-functional-A-inverse}\\
  (-\partial_\tau^2+\omega^2)G(\tau,\tau')
  &=\delta(\tau-\tau').
  \label{eq:proposal-green-equation}
\end{align}

% YIN-SOURCE: id=PROPOSAL-A-U079; arguments=YIN253A-C02-A07; notes=253a:26-27; pdf=37-38; video=M0py5a4RWhE:01:15:34.260-01:23:04.739; transcript=YIN253A-C02-T000442..YIN253A-C02-T000455; equations=YIN253A-C02-EQ011; class=SOURCE_COMPOSITE
Set $\hbar=1$ for the following normalized Gaussian functional integral.
Functional integration by parts determines its two-point function without
choosing a normalization for each discretized mode.  The derivative of the
Gaussian supplies $-Ag$, which is removed by $-A^{-1}$.

% YIN-SOURCE: id=PROPOSAL-A-U080; arguments=YIN253A-C02-A07; notes=253a:26-27; pdf=37-38; video=M0py5a4RWhE:01:15:34.260-01:23:04.739; transcript=YIN253A-C02-T000442..YIN253A-C02-T000455; equations=YIN253A-C02-EQ011; class=SOURCE_COMPOSITE
\begin{align}
  \langle g(\tau_1)g(\tau_2)\rangle
  &:=\frac{
    \displaystyle\int[Dg]\,
    e^{-\frac12(g,Ag)}g(\tau_1)g(\tau_2)}
  {\displaystyle\int[Dg]\,e^{-\frac12(g,Ag)}},
  \label{eq:proposal-functional-expectation}\\
  \int[Dg]e^{-\frac12(g,Ag)}g(\tau_1)g(\tau_2)
  &=\int[Dg]g(\tau_2)
  \left(-A^{-1}\cdot\frac{\delta}{\delta g(\tau_1)}\right)
  e^{-\frac12(g,Ag)}
  \nonumber\\
  &=\int[Dg]e^{-\frac12(g,Ag)}
  A^{-1}\cdot\frac{\delta g(\tau_2)}{\delta g(\tau_1)}
  \nonumber\\
  &=G(\tau_1,\tau_2)
  \int[Dg]e^{-\frac12(g,Ag)}.
  \label{eq:proposal-functional-integration-by-parts}
\end{align}

% YIN-SOURCE: id=PROPOSAL-A-U081; arguments=YIN253A-C02-A07; notes=253a:27; pdf=38; video=M0py5a4RWhE:01:18:31.260-01:23:04.739; transcript=YIN253A-C02-T000449..YIN253A-C02-T000455; equations=YIN253A-C02-EQ011; class=SOURCE_COMPOSITE
The inverse-kernel identity used in the last line is

% YIN-SOURCE: id=PROPOSAL-A-U082; arguments=YIN253A-C02-A07; notes=253a:27; pdf=38; video=M0py5a4RWhE:01:18:31.260-01:23:04.739; transcript=YIN253A-C02-T000449..YIN253A-C02-T000455; equations=YIN253A-C02-EQ011; class=SOURCE_COMPOSITE
\begin{equation}
  A^{-1}\cdot\delta(\tau_1-\tau_2)
  =\int d\tau'\,G(\tau_1,\tau')
  \delta(\tau'-\tau_2)
  =G(\tau_1,\tau_2),
  \qquad
  \langle g(\tau_1)g(\tau_2)\rangle=G(\tau_1,\tau_2).
  \label{eq:proposal-functional-covariance}
\end{equation}

% SOURCE-CONFLICT: note p.26 writes Dg in the numerator and Dx in the
% denominator. The proposal declares one normalized [Dg] measure. Local forms
% A.g, (Ag), and Ag are represented by the single operator convention above.

% YIN-SOURCE: id=PROPOSAL-A-U083; arguments=YIN253A-C02-A08; notes=253a:28-29; pdf=39-40; video=vk_RlYUKUyM:00:00:12.660-00:14:04.040; transcript=YIN253A-C02-T000461..YIN253A-C02-T000484; equations=YIN253A-C02-EQ012; class=SOURCE_COMPOSITE
\subsection{The Green function in Fourier space}

% YIN-SOURCE: id=PROPOSAL-A-U084; arguments=YIN253A-C02-A08; notes=253a:28; pdf=39; video=vk_RlYUKUyM:00:00:12.660-00:03:31.519; transcript=YIN253A-C02-T000461..YIN253A-C02-T000465; equations=YIN253A-C02-EQ012; class=SOURCE_COMPOSITE
Translation invariance makes the Fourier basis especially convenient.  For a
real function $g(\tau)$, the negative-frequency coefficient is the complex
conjugate of the positive-frequency coefficient.  The operator $A$ acts by
multiplication in this basis.

% YIN-SOURCE: id=PROPOSAL-A-U085; arguments=YIN253A-C02-A08; notes=253a:28; pdf=39; video=vk_RlYUKUyM:00:00:12.660-00:03:31.519; transcript=YIN253A-C02-T000461..YIN253A-C02-T000465; equations=YIN253A-C02-EQ012; class=SOURCE_COMPOSITE
\begin{align}
  g(\tau)
  &=\int_{-\infty}^{\infty}\frac{dk}{2\pi}\,
  \widetilde g(k)e^{ik\tau},
  &\widetilde g^*(k)&=\widetilde g(-k),
  \label{eq:proposal-fourier-convention}\\
  \widetilde A\cdot\widetilde g(k)
  &=(k^2+\omega^2)\widetilde g(k),
  &\widetilde A^{-1}\cdot\widetilde g(k)
  &=\frac{1}{k^2+\omega^2}\widetilde g(k).
  \label{eq:proposal-fourier-operator}
\end{align}

% YIN-SOURCE: id=PROPOSAL-A-U086; arguments=YIN253A-C02-A08; notes=253a:28; pdf=39; video=vk_RlYUKUyM:00:02:11.760-00:03:31.519; transcript=YIN253A-C02-T000463..YIN253A-C02-T000465; equations=YIN253A-C02-EQ012; class=SOURCE_COMPOSITE
Applying $-\partial_\tau^2+\omega^2$ to the inverse Fourier transform cancels
its denominator and leaves the Fourier representation of the delta function.
This verifies the Green equation directly.

% YIN-SOURCE: id=PROPOSAL-A-U087; arguments=YIN253A-C02-A08; notes=253a:28; pdf=39; video=vk_RlYUKUyM:00:02:11.760-00:03:31.519; transcript=YIN253A-C02-T000463..YIN253A-C02-T000465; equations=YIN253A-C02-EQ012; class=SOURCE_COMPOSITE
\begin{align}
  G(\tau,\tau')
  &=\int_{-\infty}^{\infty}\frac{dk}{2\pi}\,
  \frac{e^{ik(\tau-\tau')}}{k^2+\omega^2},
  \label{eq:proposal-fourier-green-function}\\
  (-\partial_\tau^2+\omega^2)G(\tau,\tau')
  &=\int_{-\infty}^{\infty}\frac{dk}{2\pi}
  e^{ik(\tau-\tau')}
  =\delta(\tau-\tau').
  \label{eq:proposal-fourier-green-check}
\end{align}

% SOURCE-CONFLICT: the gray residue line on note p.28 ends in |tau_2| even
% though the Green function uses tau and tau'. The general kernel arguments
% govern the proposal.

% YIN-SOURCE: id=PROPOSAL-A-U088; arguments=YIN253A-C02-A08; notes=253a:28-29; pdf=39-40; video=vk_RlYUKUyM:00:02:11.760-00:11:03.720; transcript=YIN253A-C02-T000463..YIN253A-C02-T000475; equations=YIN253A-C02-EQ012; class=SOURCE_COMPOSITE
If you find this confusing, always go back to the definition and see what is
going on.  The residue calculation already performed gives the same
coordinate-space kernel, $G(\tau,\tau')=(2\omega)^{-1}
e^{-\omega|\tau-\tau'|}$ in the present $\hbar=1$ convention.

% YIN-SOURCE: id=PROPOSAL-A-U089; arguments=YIN253A-C02-A08; notes=253a:29; pdf=40; video=vk_RlYUKUyM:00:11:34.459-00:14:04.040; transcript=YIN253A-C02-T000477..YIN253A-C02-T000484; equations=YIN253A-C02-EQ012; class=SOURCE_COMPOSITE
The quadratic action is also diagonal in Fourier space.  Its normalized mode
covariance carries a delta function imposing $k_1+k_2=0$, as required by the
reality condition and translation invariance.

% YIN-SOURCE: id=PROPOSAL-A-U090; arguments=YIN253A-C02-A08; notes=253a:29; pdf=40; video=vk_RlYUKUyM:00:11:34.459-00:14:04.040; transcript=YIN253A-C02-T000477..YIN253A-C02-T000484; equations=YIN253A-C02-EQ012; class=SOURCE_COMPOSITE
\begin{align}
  S^E
  &=\frac12\int_{-\infty}^{\infty}\frac{dk}{2\pi}\,
  \widetilde g(-k)(k^2+\omega^2)\widetilde g(k),
  \label{eq:proposal-fourier-action}\\
  \left\langle\widetilde g(k_1)\widetilde g(k_2)\right\rangle
  &=\frac{2\pi\delta(k_1+k_2)}{k_1^2+\omega^2}.
  \label{eq:proposal-fourier-covariance}
\end{align}

% SOURCE-CONFLICT: note p.29 calls this the \widetilde x(k) basis while its
% equations use \widetilde g(k). The equation variable is retained.

% YIN-SOURCE: id=PROPOSAL-A-U091; arguments=YIN253A-C02-A08; notes=253a:28-29; pdf=39-40; video=vk_RlYUKUyM:00:10:33.800-00:12:42.800; transcript=YIN253A-C02-T000475..YIN253A-C02-T000479; equations=YIN253A-C02-EQ012; class=SOURCE_COMPOSITE
If you want to do something different, you have to show that it works.  The
regulated mode definition, its inverse kernel, and the contour prescription
give one consistent construction.  The denominator $k^2+\omega^2$ will serve
as the free propagator when interactions are introduced.

% PROPOSAL-COVERAGE: arguments=YIN253A-C02-A01,YIN253A-C02-A02,YIN253A-C02-A03,YIN253A-C02-A04,YIN253A-C02-A05,YIN253A-C02-A06,YIN253A-C02-A07,YIN253A-C02-A08; equations=YIN253A-C02-EQ001..YIN253A-C02-EQ012; notes=253a:10-29; pdf=20-40.
% PROPOSAL-VOICE: retained or lightly recast YIN253A-C02-VOICE-0001,0003,0005,0013,0016,0018,0021,0025,0031,0035,0038,0042,0043,0046,0050,0054,0058,0059,0064,0066,0067.
% PROPOSAL-CONFLICTS-WITHHELD: undefined [Dq]; H argument order; O(hbar) glyph; p/P and N/tilde-N measure marks; barred energies; p.17 Euclidean sign and missing hat/dt; p.18 doubled-dot derivative; p.26 Dg/Dx; undefined tau_12 reconciled by A06 requirement; p.28 |tau_2| residue; p.29 x/g Fourier-basis label.
